% SEGMENT OLP-0430-S01
% Part: normal-modal-logic
% Chapter: axioms-systems
% Section: logics-proofs

\documentclass[../../../include/open-logic-section]{subfiles}

\begin{document}

\olfileid{nml}{prf}{prf}

\olsection{\usetoken{P}{derivation} மற்றும் வாய்ப்புநிலை அமைப்புகள்}

இயல்பான வாய்ப்புநிலைத் தருக்கங்களுக்கு ஒரு !!a{derivation} என்பது என்னவென்று
முதலில் வரையறுக்கிறோம். தோராயமாக, !!a{derivation} என்பது !!{formula}s இன்
ஒரு தொடர்; அதிலுள்ள ஒவ்வொரு !!{element} உம் சில \emph{அடிகோள்களில்}
ஒன்றாகவோ (அதன் ஒரு பதிலீட்டு நிகழ்வாகவோ), அல்லது முந்தைய !!{element}s
இலிருந்து சில அனுமான விதிகளில் ஒன்றால் பின்வருவதாகவோ இருக்கும். இயல்பான
வாய்ப்புநிலைத் தருக்கங்களுக்கு, மெய்மங்களின் எல்லா நிகழ்வுகளும்
\iftag{prvDiamond}{, \Ax{K}, மற்றும்~\Dual{}}{ மற்றும்~\Ax{K}}
அடிகோள்களாக எண்ணப்படுகின்றன. இதனால் மிகச்சிறிய இயல்பான வாய்ப்புநிலைத்
தருக்கமான வாய்ப்புநிலை அமைப்பு~$\Log{K}$ கிடைக்கிறது. எனினும், மற்ற
அமைப்புகளைப் பெற நாம் கூடுதல் அடிகோள்களைச் சேர்க்க விரும்பலாம். விதிகள்
எப்போதும் மோடஸ் போனென்ஸ்~\MP{} மற்றும் இன்றியமையாக்கல்~\Nec{} ஆகும்.

\begin{defn}
  ஒரு வாய்ப்புநிலை அமைப்பு~$\Log{K} !A_1 \dots !A_n$ உம் ஒரு
  !!a{formula}~$!B$ உம் கொடுக்கப்பட்டுள்ளன என்க. $!C_k = !B$ ஆகவும்,
  ஒவ்வொரு~$!C_i$ உம் ஒரு மெய்ம நிகழ்வாகவோ, $\Ax{K}$,
  \iftag{prvDiamond}{ $\Dual$,}{} $!A_1$, \dots,~$!A_n$ ஆகியவற்றில்
  ஒன்றின் நிகழ்வாகவோ, அல்லது முந்தைய !!{formula}s இலிருந்து~\MP{} அல்லது
  \Nec{} விதிகளால் பின்வருவதாகவோ இருக்கும் !!{formula}s~$!C_1$,
  \dots,~$!C_k$ இருந்தால், இருந்தால் மட்டுமே, $!B$ ஆனது
  $\Log{K} !A_1 \dots !A_n$ இல் \emph{!!{derivable}} என்கிறோம்; இதனை
  $\Log{K} !A_1 \dots !A_n \Proves !B$ என்று எழுதுகிறோம்.
\end{defn}

பின்வரும் கூற்று,~$!B$ இன் ஒரு~$\Sigma$-!!{derivation} ஐக் காட்டுவதன்
மூலம் $!B \in \Sigma$ என்று காட்ட அனுமதிக்கிறது.

\begin{prop}
  $\Log{K} !A_1 \dots !A_n = \Setabs{!B}{\Log{K} !A_1 \dots !A_n \Proves !B}$.
\end{prop}

\begin{proof}
  பின்வருவதைக் காட்ட !!{derivation}s இன் நீளம் மீதான தொகுத்தறிதலைப்
  பயன்படுத்துகிறோம்:
  $\Setabs{!B}{\Log{K} !A_1 \dots !A_n \Proves !B} \subseteq
  \Log{K} !A_1 \dots !A_n$.

  $!B$ இன் !!{derivation} இன் நீளம்~$1$ எனில், அது ஒரே ஒரு
  !!a{formula} ஐக் கொண்டுள்ளது. அந்த !!a{formula} முந்தைய வாய்பாடுகளிலிருந்து
  \MP{} அல்லது~\Nec{} ஆல் பின்வர முடியாது; ஆகவே அது ஒரு மெய்ம நிகழ்வாக,
  \Ax{K} இன் நிகழ்வாக,\iftag{prvDiamond}{ $\Dual$ இன் நிகழ்வாக,}{}
  அல்லது $!A_1$, \dots,~$!A_n$ ஆகியவற்றில் ஒன்றின் நிகழ்வாக இருக்க வேண்டும்.
  ஆனால் $\Log{K}!A_1\dots!A_n$ உம் இவற்றைக் கொண்டுள்ளது; ஆகவே
  $!B \in \Log{K}!A_1 \dots !A_n$.

  $!B$ இன் !!{derivation} இன் நீளம்~$> 1$ எனில், அதோடு $!B$ ஆனது
  !!{derivation} இன் கடைசி வரியாக வராத !!{formula}s இலிருந்து~\MP{}
  அல்லது~\Nec{} ஆல் பெறப்பட்டிருக்கலாம். $!B$ ஆனது $!C$, $!C \lif !B$
  ஆகியவற்றிலிருந்து (\MP{} ஆல்) பின்வந்தால், தொகுத்தறிதல் கருதுகோளின்படி
  $!C$, $!C \lif !B \in \Log{K}!A_1 \dots !A_n$. ஆனால் ஒவ்வொரு
  வாய்ப்புநிலைத் தருக்கமும் மோடஸ் போனென்ஸின் கீழ் மூடப்பட்டுள்ளது; ஆகவே
  $!B \in \Log{K}!A_1 \dots !A_n$. $!B \equiv \Box !C$ ஆனது $!C$
  இலிருந்து~\Nec{} ஆல் பின்வந்தால், தொகுத்தறிதல் கருதுகோளின்படி
  $!C \in \Log{K}!A_1 \dots !A_n$. ஆனால் ஒவ்வொரு இயல்பான வாய்ப்புநிலைத்
  தருக்கமும்~$\Nec$ இன் கீழ் மூடப்பட்டுள்ளது; ஆகவே
  $!B \in \Log{K}!A_1\dots!A_n$.

  மறுதலை உட்கிடைப்பைக் காட்ட,~$\Sigma = \Setabs{!B}{\Log{K} !A_1 \dots
  !A_n \Proves !B }$ ஆனது $!A_1$, \dots,~$!A_n$ ஆகியவற்றின் எல்லா
  நிகழ்வுகளையும் கொண்ட ஓர் இயல்பான வாய்ப்புநிலைத் தருக்கம் என்று காட்டுகிறோம்;
  மேலும் $\Log{K} !A_1 \dots !A_n$ என்பது, வரையறைப்படி, அத்தகைய
  மிகச்சிறிய தருக்கம் என்ற கூற்றைப் பயன்படுத்துகிறோம்.
  \begin{enumerate}
    \item ஒவ்வொரு மெய்மம்~$!B$ உம் ஒரு மெய்ம நிகழ்வு; ஆகவே
      $\Log{K}!A_1\dots!A_n \Proves !B$; எனவே $\Sigma$ எல்லா
      மெய்மங்களையும் கொண்டுள்ளது.
    \item $\Log{K}!A_1\dots!A_n \Proves !C$ மற்றும்
      $\Log{K}!A_1\dots!A_n \Proves !C \lif !B$ எனில்,
      $\Log{K}!A_1\dots!A_n \Proves !B$: $!C$ இன் !!{derivation} உடன்
      $!C \lif !B$ க்கானதையும் இணைத்து,~$!B$ என்ற வரியைச்
      சேர்க்க. கடைசி வரி~\MP{} ஆல் நியாயப்படுத்தப்படுகிறது. ஆகவே~$\Sigma$
      மோடஸ் போனென்ஸின் கீழ் மூடப்பட்டுள்ளது.
    \item $!B$ க்கு ஓர் !!a{derivation} இருந்தால்,~$!B$ இன் ஒவ்வொரு
      பதிலீட்டு நிகழ்வுக்கும் ஒரு !!a{derivation} உண்டு: அந்தப் பதிலீட்டை
      !!{derivation} இலுள்ள ஒவ்வொரு !!a{formula} க்கும் பயன்படுத்துக.
      (பயிற்சி: !!{derivation}s இன் நீளம் மீதான தொகுத்தறிதலால் கிடைக்கும்
      விளைவும் சரியான !!a{derivation} தான் என்று நிறுவுக.) ஆகவே~$\Sigma$
      சீரான பதிலீட்டின் கீழ் மூடப்பட்டுள்ளது. ($\Sigma$ ஒரு வாய்ப்புநிலைத்
      தருக்கத்தின் எல்லா நிபந்தனைகளையும் நிறைவுசெய்கிறது என்று இப்போது
      நிறுவிவிட்டோம்.)
    % SOURCE CORRECTION TA-NML-016: membership is of the axiom schema \Ax{K}, not an unrelated formula K.
    \item $\Log{K}!A_1\dots!A_n \Proves \Ax{K}$ நமக்கு உள்ளது; ஆகவே
      $\Ax{K} \in \Sigma$.
    \tagitem{prvDiamond}{$\Log{K}!A_1\dots!A_n \Proves \Dual$ நமக்கு
      உள்ளது; ஆகவே $\Dual \in \Sigma$.}{}
    \item $\Log{K}!A_1\dots!A_n \Proves !C$ எனில், கூடுதல்
      வரி~$\Box !C$ ஆனது~\Nec{} ஆல் நியாயப்படுத்தப்படுகிறது. இதன் விளைவாக,
      $\Sigma$ ஆனது~\Nec{} இன் கீழ் மூடப்பட்டுள்ளது. ஆகவே~$\Sigma$
      இயல்பானது.
  \end{enumerate}
\end{proof}

\end{document}
